\newif\ifcomments     %% include author discussion
\newif\ifanonymous    %% include author identities
\newif\ifextended     %% include appendix
\newif\ifsubmission   %% prepare the submitted version
\newif\ifpublic       %% version available for posting / final version

\commentstrue         %% toggle comments here
\extendedfalse

\submissionfalse      %% at most one of these must be true (neither for draft version)
\publicfalse          %% but if you want to see comments, these should both be off

% If we are going to make a version public, i.e. on arXiv, we should 
% make sure that there are no comments and our names are on it.
\ifpublic
\submissionfalse
\commentsfalse
\anonymousfalse
\fi

%% If we are submission, make sure there are no comments and our names 
%% are NOT on it.
\ifsubmission
\publicfalse
\commentsfalse
\anonymoustrue
\fi

\documentclass[acmsmall, screen=true, nonacm=true,
  review=\ifpublic false\else true\fi,
  anonymous=\ifanonymous true\else false\fi
]{acmart}

% \setcopyright{rightsretained}
% \acmDOI{}
% \acmYear{2025}
% \copyrightyear{2025}
% \acmSubmissionID{}
% \acmJournal{PACMPL}
% \acmVolume{}
% \acmNumber{POPL}
% \acmArticle{}
% \acmMonth{}
% \received{}
% \received[accepted]{}

\let\Bbbk\relax
\usepackage{draft} % local
\usepackage[supertabular]{ottalt} % local
\usepackage{amsmath,amssymb}
\usepackage{natbib}
\usepackage{mathpartir,mathtools,stmaryrd}
\usepackage{listings,xspace,color,doi}
\usepackage[para,flushmargin,multiple]{footmisc}
\usepackage[noend]{algpseudocode}
\usepackage[capitalize,noabbrev,nameinlink]{cleveref}

\citestyle{acmauthoryear}

\hypersetup{
  colorlinks=true,
  urlcolor=blue,
  citecolor=violet,
  backref=true
}

\lstset{
  basicstyle=\sffamily,
  columns=flexible,
  mathescape=true,
  breaklines=false,
  morekeywords={data,type,case,of,if,then,else,let,in,refl,transp,J},
  %stringstyle=\color{string},
  %morestring=*[d]{"}, % chktex 18
  %morestring=*[d]{'},
  literate=
    {->}{{$\mathrel{\rightarrow}$}}2
    {\\}{{$\lambda$}}1
    {=>}{{$\mathrel{\Rightarrow}$}}2
    {|-}{{$\vdash$}}1
    {~}{{$\mathbin{\sim}$}}1
    {'}{{$^\prime$}}0
}
\def\lstinlinecol{\lstinline[columns=fullflexible]}

\ifcomments
\newnote{sw}{blue} % Stephanie Weirich
\newnote{yl}{purple} % Yiyun Liu
\newnote{jc}{olive} % Jonathan Chan
\else
\newcommand{\sw}[1]{}
\newcommand{\yl}[1]{}
\newcommand{\jc}[1]{}
\fi

\newcommand{\suppress}[1]{}
\newcommand{\name}{DCOI\@\xspace}
\newcommand{\longname}{Dependent Calculus of Indistinguishability\@\xspace}
\newcommand{\subst}[3]{#1[#2 \mapsto #3]}
\newcommand{\dotv}[3]{\href{#1}{\texttt{#2}}{\texttt{:#3}}}
\newcommand{\dotpi}[2]{\href{#1/#2}{\texttt{#2}}}

% e.g. \punctstack{.}\footnote{...}
% stacks the footnote marker above .
\newlength{\punctwidth}
\newcommand{\punctstack}[1]{#1%
  \settowidth{\punctwidth}{#1}%
  \hspace*{-\the\punctwidth}%
}

\inputott{dcoi-rules}

\title{DCOI}

\author{Yiyun Liu}
\orcid{0009-0006-8717-2498}
\affiliation{%
  \institution{University of Pennsylvania}
  \city{Philadelphia}
  \country{USA}
}
\email{liuyiyun@seas.upenn.edu}

\author{Jonathan Chan}
\orcid{0000-0003-0830-3180}
\affiliation{%
  \institution{University of Pennsylvania}
  \city{Philadelphia}
  \country{USA}
}
\email{jcxz@seas.upenn.edu}

\author{Stephanie Weirich}
\orcid{0000-0002-6756-9168}
\affiliation{%
  \institution{University of Pennsylvania}
  \city{Philadelphia}
  \country{USA}
}
\email{sweirich@seas.upenn.edu}

%% By default, the full list of authors will be used in the page
%% headers. Often, this list is too long, and will overlap
%% other information printed in the page headers. This command allows
%% the author to define a more concise list
%% of authors' names for this purpose.
% \renewcommand{\shortauthors}{Liu et al.}

\begin{CCSXML}
<ccs2012>
  <concept>
    <concept_id>10003752.10003790.10011740</concept_id>
    <concept_desc>Theory of computation~Type theory</concept_desc>
    <concept_significance>500</concept_significance>
  </concept>
</ccs2012>
\end{CCSXML}
\ccsdesc[500]{Theory of computation~Type theory}
\keywords{Modes, Dependent Types, Coq, Formalization}

\begin{document}

\begin{abstract}

\end{abstract}

\maketitle

\begin{spreadlines}{0em}
\begin{align*}
  [[:concrete: Ifz & : H ! Nat -> H Univ 0 -> H Univ 0 -> Univ 0 @ H]] \\
  [[:concrete: Ifz & := \ H n . \ H A . \ H B . phantom]] \\
  & [[:concrete: case n return Univ 0 of
      & Z => A ++
      & S n => B]] \\
  [[:concrete: discr & : Pi A : H Univ 0 . Pi n : H ! Nat . H (! Z ~ H ! S H n : ! Nat) -> A @ H]] \\
  [[:concrete: discr & := \ H A . \ H n . \ H p . J ! Z p]] \\ % (\ H m . \ H q . ! Ifz H m H ! Nat H A)
  [[:concrete: head & : Pi A : H Univ 0 . Pi n : H ! Nat . L ! Vec H A H n -> A @ L]] \\
  [[:concrete: head & := \ H A . \ H n . \ L v . phantom]] \\
  & [[:concrete: (case v return \ H m . L (m ~ H ! S H n : ! Nat) -> A of
      & vnil => \ H p . absurd (! discr H Void H n H p) ++
      & vcons n x xs => \ H p . x) H refl]]
\end{align*}
\end{spreadlines}

\begin{alignat*}{2}
  [[G]] &\Coloneqq [[nil]] \mid [[G ++ x : l A]] \\
  [[E]] &\Coloneqq [[nil]] \mid [[E ++ x : l]]
\end{alignat*}
\nonterms{tm}
\drules[Wt]{$[[G |- a : l A]]$}{Typing rules}{Var,Pi,Abs,App,Univ,Void,Absurd,Eq,Refl,J,Sig,Pack,Let,Conv}
\drules[Wtg]{$[[G |- a : A]]$}{Existentially-graded typing}{Grade}
\drules[Wf]{$[[ |- G]]$}{Context well-formedness rules}{Nil, Cons}
\drules[I]{$[[E |- a == b @ l]]$}{Indistinguishability rules}{Var,Pi,Abs,App,Univ,Void,Absurd,Eq,Refl,J,Sig,Pack,Let}
\drules[GI]{$[[E |- a == b @ l l0]]$}{Guarded indistinguishability rules}{Dist,Indist}
\drules[E]{$[[E |- a == b]]$}{Definitional equality}{Conv}
\drules[Wg]{$[[E |- a : l]]$}{Grading rules}{Var,Pi,Abs,App,Univ,Void,Absurd,Eq,Refl,J,Sig,Pack,Let}

\cite{liu:dcoi}

\begin{acks} 
The authors would like to thank the anonymous reviewers for their comments and suggestions.
This work was supported by the \grantsponsor{}{National Science Foundation}{} under Grant Nos.~\grantnum{}{2006535} and \grantnum{}{2327738}.
\end{acks}

\bibliographystyle{ACM-Reference-Format}
\bibliography{refs}

\end{document}